\begin{definition}[Sup-convolution, (3.13) and (3.14)]
Let $\mathbb G$ be a Carnot setting (\noderef{CarnotSetting}) with group law $*$ and
homogeneous norm $\|\cdot\|$, put $q:=2\,r!$, let $\Omega\subset\mathbb R^n$, let
$w:\mathbb R^n\to\mathbb R$, let $\varepsilon>0$ and let $x\in\mathbb R^n$. The
\emph{sup-convolution} of $w$ is
\[
 w^{\varepsilon}(x)\ :=\ \sup\Big\{\,w(z*x)-\frac{\|z\|^{q}}{2\varepsilon}\ :\
   z\in\mathbb R^n,\ z*x\in\overline\Omega\,\Big\} .
\]

This is (3.14), which is (3.13),
$w^{\varepsilon}(x)=\sup_{y\in\overline\Omega}\{w(y)-\|y*x^{-1}\|^{q}/2\varepsilon\}$,
after the substitution $y=z*x$: the sup-convolution is a supremum of \emph{left}
translates of $w$, which is what makes the kernel compatible with the equation, since the
horizontal vector fields are left-invariant. The exponent $q=2\,r!$ is chosen so that
$\|\cdot\|^{q}$ is a polynomial, hence smooth; this is the kernel used by Wang for Carnot
groups.

The supremum is taken over the set of admissible $z$, which is non-empty exactly when
$x\in\overline\Omega$, since $z=0$ gives $0*x=x$ and $\|0\|=0$ --- in which case
$w^{\varepsilon}(x)\ge w(x)$. It is finite whenever $\Omega$ is bounded and $w$ is
continuous on $\overline\Omega$, because then $w$ is bounded on the compact set
$\overline\Omega$ and the penalty $-\|z\|^{q}/2\varepsilon$ is non-positive, so the family
is bounded above by $\sup_{\overline\Omega}w$. In the formal statement the supremum is the
supremum of the indexed family of real numbers, which by convention returns $0$ when the
family is empty or unbounded above; the nodes that use $w^{\varepsilon}$ therefore carry
the hypotheses $\Omega$ bounded, $w\in C(\overline\Omega)$ and $x\in\overline\Omega$ under
which neither degenerate case occurs, and under which the above is the supremum in the
ordinary sense.
\end{definition}
